\chapter{Experimental Design}
\label{chap:design}

The derivation in Chapter~\ref{chap:expanded} establishes a plausible
finite-sample reference but not an exact distributional theorem. Its value
must therefore be assessed through repeated sampling. This chapter describes
one controlled experiment designed to answer a direct question: when two
different categorical populations have exactly equal positive MI, how often
does each analytic method incorrectly reject equality?

\section{Validation principle}

For each scenario, \(P\) and \(Q\) are fixed \(r\times c\) population
probability tables satisfying
\begin{equation}
  I(P)=I(Q).
  \label{eq:design-equal-mi}
\end{equation}
Independent count tables are repeatedly drawn as
\begin{equation}
  N^{P}\sim\operatorname{Multinomial}(n_P,P),
  \qquad
  N^{Q}\sim\operatorname{Multinomial}(n_Q,Q).
  \label{eq:design-multinomial}
\end{equation}
Normal Wald, simple Welch, and expanded Welch analyse the same sampled table
pair in every replicate. This pairing removes Monte Carlo variation from
between-method differences and isolates the effect of the reference
distribution.

The primary significance level is \(\alpha=0.05\). A calibrated method should
reject Equation~\eqref{eq:design-equal-mi} in approximately 5\% of repeated
samples. Additional checks use \(\alpha=0.10\), \(\alpha=0.01\), and a full
rejection-calibration curve from 0 to 0.10.

\section{Constructing distinct populations with equal MI}

The weak null permits \(P\ne Q\). Each scenario therefore requires two
different joint distributions with numerically equal MI. The construction
separates marginal shape from association strength.

\subsection{Random margins}

For a variable with \(m\) categories, a symmetric Dirichlet distribution
generates a positive probability vector
\begin{equation}
  (u_1,\ldots,u_m)
  \sim\operatorname{Dirichlet}(a,\ldots,a).
  \label{eq:dirichlet-margin}
\end{equation}
The concentration \(a\) controls how similar the category probabilities are.
Large values produce margins close to uniform. Values below one commonly
produce several rare categories. The implementation rejects margins with a
probability below \(10^{-4}\) or above 0.98 to avoid numerical degeneracy while
retaining substantial skewness.

The row and column margins for \(P\) and \(Q\) are generated independently.
The two populations can therefore have different marginal distributions even
when their MI values are equal.

\subsection{Random association pattern}

Margins alone define the independent table. To introduce association, the
generator draws a matrix \(A=(A_{ij})\) of independent standard-normal values,
removes its row and column means, and standardises its root-mean-square
magnitude. This produces a non-additive interaction pattern that cannot be
absorbed into the margins.

For association strength \(\lambda\), define a positive log-linear kernel
\begin{equation}
  K_{ij}(\lambda)=\exp(\lambda A_{ij}).
\end{equation}
Iterative proportional fitting finds positive scaling vectors \(u\) and \(v\)
such that
\begin{equation}
  p_{ij}(\lambda)=u_iK_{ij}(\lambda)v_j
  \label{eq:loglinear-table}
\end{equation}
has the required row and column margins. A one-dimensional root finder then
selects \(\lambda\) so that
\begin{equation}
  I\{P(\lambda)\}=I_{\mathrm{target}}.
\end{equation}
The same procedure is run independently for \(Q\). Its margins, interaction
pattern, and fitted association strength differ from those of \(P\), but its
target MI is identical.

A generated pair is retained only when all cells have positive population
probability, both first-order MI variances are positive, and
\begin{equation}
  \lVert P-Q\rVert_1
  =\sum_{i,j}|p_{ij}-q_{ij}|
  \ge 0.05.
  \label{eq:l1-population-separation}
\end{equation}
The largest absolute difference between \(I(P)\) and \(I(Q)\) among the 60
accepted pairs was \(8.5\times10^{-14}\) nats. Numerical mismatch in the null
was therefore negligible relative to sampling variation.

\section{Table shapes and regimes}

Every regime uses six table shapes:
\begin{equation}
  2\times2,\quad 2\times5,\quad 3\times3,\quad
  3\times5,\quad 5\times5,\quad 8\times8.
  \label{eq:design-shapes}
\end{equation}
Each regime has two variants. The complete design contains
\begin{equation}
  6\text{ shapes}\times5\text{ regimes}\times2\text{ variants}
  =60
\end{equation}
fixed equal-MI population pairs. Both sample sizes lie between 50 and 1,000.
The two target values, 0.10 and 0.15 nats, keep the primary experiment away
from the nonregular independence boundary.

Table~\ref{tab:regime-designs} gives the exact construction rules. The
concentrations \((a_P,a_Q)\) apply to both row and column margins in the
corresponding populations. The ratio is \(n_Q:n_P\).

\begin{table}[htbp]
\centering
\caption{Exact equal-MI regime specifications.}
\label{tab:regime-designs}
\small
\begin{tabularx}{\textwidth}{@{}XcccX@{}}
\toprule
Regime and variant & MI & \((a_P,a_Q)\) & Ratio & Sampling rule \\
\midrule
Well sampled 1 & 0.10 & \((50,50)\) & \(1:1\) & 15 observations per cell; \(n_P\ge200\) \\
Well sampled 2 & 0.15 & \((50,50)\) & \(1:1\) & 15 observations per cell; \(n_P\ge200\) \\
Moderate 1 & 0.10 & \((8,4)\) & \(1:1\) & 8 observations per cell; \(n_P\ge100\) \\
Moderate 2 & 0.15 & \((8,2)\) & \(2:1\) & 6 observations per cell; \(n_P\ge100\) \\
Highly skewed and sparse 1 & 0.10 & \((2,2)\) & \(1:1\) & Both minima in \([1,5)\) \\
Highly skewed and sparse 2 & 0.15 & \((10,2)\) & \(2:1\) & Both minima in \([1,5)\) \\
Ultra-skewed and sparse 1 & 0.10 & \((0.8,0.8)\) & \(1:1\) & Both minima in \([0.05,1)\) \\
Ultra-skewed and sparse 2 & 0.15 & \((10,0.8)\) & \(10:1\) & Both minima in \([0.05,1)\) \\
Widespread sparsity 1 & 0.10 & \((0.35,0.35)\) & \(1:1\) & 25--50\% of cells below 1 and at least 50\% below 5 \\
Widespread sparsity 2 & 0.15 & \((1,0.35)\) & \(2:1\) & Same cell-wide sparsity constraints \\
\bottomrule
\end{tabularx}
\end{table}

For the well-sampled and moderate regimes, sample size is fixed from the shape
and design constants:
\begin{equation}
  n_P=\max(n_{\min},drc),
  \qquad
  n_Q=\rho n_P,
  \label{eq:fixed-density-n}
\end{equation}
where \(d\) is the planned average number of observations per cell and
\(\rho=n_Q/n_P\). This rule is chosen before the population tables are drawn.

For the three sparse regimes, sample sizes are selected after the population
tables are constructed. The selection uses true joint expected counts
\begin{equation}
  E^{P}_{ij}=n_Pp_{ij},
  \qquad
  E^{Q}_{ij}=n_Qq_{ij}
  \label{eq:true-expected-counts}
\end{equation}
to meet the stated rare-cell or cell-wide sparsity constraint. The resulting
sample sizes and smallest expected counts are summarised in
Table~\ref{tab:realized-designs}.

\begin{table}[htbp]
\centering
\caption{Realised sample-size and expected-count ranges across the six shapes.}
\label{tab:realized-designs}
\small
\begin{tabularx}{\textwidth}{@{}Xccc@{}}
\toprule
Regime and variant & \(n_P\) & \(n_Q\) & Smallest expected count \\
\midrule
Well sampled 1 & 200--960 & 200--960 & 2.140--18.228 \\
Well sampled 2 & 200--960 & 200--960 & 1.385--15.015 \\
Moderate 1 & 100--512 & 100--512 & 0.790--6.927 \\
Moderate 2 & 100--384 & 200--768 & 0.057--1.178 \\
Highly skewed and sparse 1 & 52--975 & 52--975 & 1.025--4.072 \\
Highly skewed and sparse 2 & 54--460 & 108--920 & 1.096--4.201 \\
Ultra-skewed and sparse 1 & 64--604 & 64--604 & 0.052--0.413 \\
Ultra-skewed and sparse 2 & 51--75 & 510--750 & 0.057--0.632 \\
Widespread sparsity 1 & 50--435 & 50--435 & below 0.001--0.042 \\
Widespread sparsity 2 & 50--414 & 100--828 & below 0.001--0.312 \\
\bottomrule
\end{tabularx}
\end{table}

Widespread sparsity is a boundary check rather than a recommended operating
regime. Its purpose is to reveal what happens when many observed cells, and
occasionally complete empirical margins, can be empty.

\section{Replication and evaluation metrics}

Each fixed population pair receives 10,000 independently sampled table pairs:
\begin{equation}
  60\times10{,}000=600{,}000
\end{equation}
null replicates, containing 1.2 million individual count tables. The
population-generation seed is 2,026,080,501 and the sampling seed is
2,026,080,502. Scenario-specific seeds are derived deterministically from
these values.

At a true rejection probability of 0.05, 10,000 replicates give Monte Carlo
standard error
\begin{equation}
  \sqrt{\frac{0.05(1-0.05)}{10{,}000}}
  \approx0.00218
\end{equation}
for one scenario. Wilson intervals are reported for scenario-level rejection
rates.

The experiment uses the following metrics.

\begin{description}[style=nextline,leftmargin=34mm,labelwidth=32mm]
  \item[False-positive rate.]
  The fraction of valid null replicates with p-value no larger than the nominal
  significance level.

  \item[Mean absolute calibration error.]
  For a regime with scenario FPRs \(\widehat f_s\),
  \begin{equation}
    \operatorname{MAE}(\alpha)
    =\frac{1}{S}\sum_{s=1}^{S}|\widehat f_s-\alpha|.
  \end{equation}
  This metric prevents liberal and conservative scenario errors from
  cancelling.

  \item[Rejection-calibration curve.]
  The empirical rejection rate is evaluated at nominal levels from 0 to 0.10
  in increments of 0.001. A calibrated test follows the diagonal.

  \item[Coverage.]
  The fraction of valid null replicates whose method-specific 95\% confidence
  interval contains the true difference, zero.

  \item[Valid rate.]
  The fraction of replicates producing a finite statistic and p-value and,
  where required, positive finite effective degrees of freedom.

  \item[Power.]
  The fraction of valid alternative replicates that correctly reject equal MI
  at \(\alpha=0.05\).
\end{description}

Regime summaries weight every fixed population pair equally. Scenario-level
results remain available so that an average improvement cannot conceal a
specific failure.

\section{Power design}

The power experiment uses five \(3\times3\) alternatives. Population \(P\)
has uniform row and column margins and \(I(P)=0.05\) nats. Population \(Q\)
has margins \((0.9,0.05,0.05)\) and a larger MI. Both use the same ordinal
interaction pattern, which keeps the comparison focused on MI difference and
sample size. Table~\ref{tab:power-design} gives the exact settings; each
receives 10,000 replicates.

\begin{table}[htbp]
\centering
\caption{Alternative-hypothesis settings for the power experiment.}
\label{tab:power-design}
\begin{tabular}{@{}ccccc@{}}
\toprule
\(I(P)\) & \(I(Q)\) & \(n_P\) & \(n_Q\) & Minimum expected counts \((P,Q)\) \\
\midrule
0.05 & 0.07 & 300 & 300 & \((18.217,1.543)\) \\
0.05 & 0.10 & 150 & 150 & \((9.108,0.926)\) \\
0.05 & 0.10 & 300 & 300 & \((18.217,1.853)\) \\
0.05 & 0.10 & 600 & 600 & \((36.433,3.705)\) \\
0.05 & 0.15 & 300 & 300 & \((18.217,1.143)\) \\
\bottomrule
\end{tabular}
\end{table}

\section{Runtime design}

Runtime is measured after count-table construction. One fixed table pair is
used for each of four shapes: \(2\times2\), \(3\times3\), \(5\times5\), and
\(8\times8\). Each complete method call is warmed up and then timed 200 times.
The median is reported. All methods use the same vectorised implementation
path, so the comparison measures the additional reference calculation rather
than differences between software packages.

\section{Reproducibility controls}

The full experiment records the fixed population tables, sample sizes,
scenario seeds, simulation seeds, method outputs, software versions, and SHA-256
hash of the experiment script. Population construction and evaluation use
separate deterministic random-number streams. No empirical result from a
replicate is used to alter its generating population or select its method.

The experiment used Python 3.13.5, NumPy 2.4.4, SciPy 1.17.1, and pandas
3.0.3. Appendix~\ref{app:reproducibility} gives the commands and output map.
